% Lösungen Einheit 1
\einheitenkopf{Einheit 1 von 4 · Wurzelziehen und Lösbarkeit}
\zweigzeile{Hier lernst du, Gleichungen wie $x^2 = c$ durch Wurzelziehen zu lösen und zu entscheiden, ob sie zwei, eine oder keine Lösung haben · neu oder schon bekannt – je nach Buch (an der Oberschule Klasse 9 oder 10) · P10 · baut auf: Quadrieren, Quadratwurzeln, lineare Gleichungen, Scheitelpunktform}
\erg{10}{a) quadratisch \quad b) linear \quad c) quadratisch \quad d) linear \quad e) quadratisch ($x\cdot x = x^2$)}
\erg{11}{a) positiv, zwei \quad b) negativ, keine \quad c) null, eine \quad d) positiv, zwei \quad e) negativ, keine}
\erg{12}{a) $x_1 = 5$, $x_2 = -5$ \quad b) $x_1 = 10$, $x_2 = -10$ \quad c) $x_1 = 12$, $x_2 = -12$ \quad d) $x_1 = 1$, $x_2 = -1$ \quad e) $x^2 = 225$, $x_1 = 15$, $x_2 = -15$; eine Seite ist $15\,\text{m}$ lang, $-15$ passt nicht, eine Länge ist nie negativ \quad f) $x_1 = \sqrt{14} \approx 3{,}74$, $x_2 = -\sqrt{14} \approx -3{,}74$ \quad g) ein Quadrat ist nie negativ, also gibt es keine Zahl mit $x^2 = -64$}
\erg{13}{a) $x^2 = 25$; $x_1 = 5$, $x_2 = -5$ \quad b) $x^2 = 49$; $x_1 = 7$, $x_2 = -7$ \quad c) $2x^2 = 32$, $x^2 = 16$; $x_1 = 4$, $x_2 = -4$ \quad d) $r^2 = \frac{V}{\pi\cdot h} = \frac{850}{\pi\cdot 12} \approx 22{,}55$; $r \approx 4{,}7\,\text{cm}$ \quad e) $4x^2 = 0$, $x^2 = 0$; eine Lösung: $x = 0$ \quad f) $2x^2 = -8$, $x^2 = -4$; keine Lösung}
\erg{14}{a) $x - 1 = 3$ oder $x - 1 = -3$; $x_1 = 4$, $x_2 = -2$ \quad b) $x_1 = 7$, $x_2 = 3$ \quad c) $x + 2 = 7$ oder $x + 2 = -7$; $x_1 = 5$, $x_2 = -9$ \quad d) $x + 6 = 0$; $x = -6$ \quad e) $(-5 + 2)^2 = (-3)^2 = 9$, $9 = 9$ (wA)}
\erg{15}{a) zwei: $x_1 = -1$, $x_2 = 3$ \quad b) eine: $x = 1$ \quad c) keine \quad d) jede Zahl kleiner als $-3$, z.\,B. $c = -5$}
\erg{16}{a) erste Aussage wahr ($x = -5$ ist die einzige Lösung); zweite falsch: $x + 6 = 3$ oder $x + 6 = -3$, also $x_1 = -3$, $x_2 = -9$ \quad b) z.\,B. $x^2 + 3 = 0$, denn $x^2 = -3$, und ein Quadrat ist nie negativ}
\erg{17}{a) Fehler: nur die positive Wurzel gezogen, die Lösung zu $x + 1 = -9$ fehlt; richtig: $x + 1 = 9$ oder $x + 1 = -9$, $x_1 = 8$, $x_2 = -10$ \quad b) $x + 3 = 8$ oder $x + 3 = -8$; $x_1 = 5$, $x_2 = -11$}
\erg{18}{a) $40$ ist positiv: $\sqrt{40}$ und $-\sqrt{40}$ ergeben quadriert beide $40$ \quad b) das Quadrat einer Zahl ist nie negativ \quad c) nein: $0^2 = 0$, also ist $x = 0$ eine Lösung – genau eine}
